\documentclass[12pt,a4paper]{report}

\usepackage[margin=1in]{geometry}
\usepackage[T1]{fontenc}
\usepackage{graphicx}
\usepackage{hyperref}
\usepackage{array}
\usepackage{tabularx}
\usepackage{xcolor}
\usepackage{float}
\usepackage{listings}
\usepackage{enumitem}

\hypersetup{
    colorlinks=true,
    linkcolor=blue,
    urlcolor=blue,
    pdftitle={OMNIRA Classroom Collaboration App Project Report}
}

\setlength{\parskip}{0.5em}
\setlength{\parindent}{0pt}
\renewcommand{\arraystretch}{1.35}

\newcommand{\reporttitle}{OMNIRA: Classroom Collaboration App}
\newcommand{\fillin}[1]{\textit{[Enter #1]}}
\newcommand{\memberrow}[1]{#1 & \fillin{team member name} & \fillin{student ID} \\ \hline}

\begin{document}

\begin{titlepage}
    \centering
    \vspace*{2cm}
    {\Huge\bfseries \reporttitle\par}
    \vspace{1.5cm}
    {\Large Project Report\par}
    \vfill
    {\large
    Course: \fillin{course name}\\[0.5em]
    Course code: \fillin{course code}\\[0.5em]
    University: \fillin{university name}\\[0.5em]
    Lecturer: \fillin{lecturer name}
    \par}
    \vfill
    {\large Submission date: \fillin{submission date}\par}
\end{titlepage}

\pagenumbering{roman}

\chapter*{Team and Course Information}
\addcontentsline{toc}{chapter}{Team and Course Information}

\textbf{Course name:} \fillin{course name}\\
\textbf{Course code:} \fillin{course code}\\
\textbf{University name:} \fillin{university name}\\
\textbf{Lecturer name:} \fillin{lecturer name}

\vspace{1em}
\textbf{Project team members}

Fill in one row for each team member and remove any unused rows.

\begin{center}
\begin{tabularx}{\textwidth}{|>{\centering\arraybackslash}p{1.2cm}|X|X|}
\hline
\textbf{No.} & \textbf{Team member name} & \textbf{Student ID} \\ \hline
\memberrow{1}
\memberrow{2}
\memberrow{3}
\memberrow{4}
\memberrow{5}
\memberrow{6}
\end{tabularx}
\end{center}

\tableofcontents
\clearpage
\pagenumbering{arabic}

\chapter{Executive Summary}
OMNIRA is an Android classroom collaboration application designed to bring common
class activities into one mobile app. It supports classroom membership, announcements,
study materials, assignments, routines, attendance, notifications, class chat, and
direct messaging. The client is built with Kotlin and Jetpack Compose and uses Firebase
Authentication and Cloud Firestore for identity and application data. Cloudinary is
used for media uploads.

This report describes the current application, its main workflows and technical
structure, and important production considerations. It also provides a dedicated
classroom-interface page for inserting a screenshot and a team information page for
course and member details.

\chapter{Introduction}
\section{Background}
Classroom activities often rely on separate channels for announcements, learning
resources, assignments, attendance, schedules, and communication. OMNIRA explores a
single mobile interface for these activities, allowing students and class administrators
to access classroom information and communication features from an Android device.

\section{Problem Statement}
When classroom communication and resources are spread across unrelated tools, users
may need to switch between apps to find updates, submit work, check schedules, or
contact classmates. A unified classroom application can make these workflows easier
to find and manage.

\section{Project Aim}
The aim of OMNIRA is to provide an Android application that combines essential
classroom information management and communication features in one user interface.

\section{Objectives}
\begin{itemize}[leftmargin=*]
    \item Provide account registration, sign-in, profile management, and logout.
    \item Allow users to create classrooms or request access using a class code.
    \item Provide class announcements, materials, assignments, routines, and attendance.
    \item Support class chat, direct conversations, and user-facing notifications.
    \item Display a connectivity warning when the device has no validated internet access.
    \item Provide a deliberate account-deletion workflow with password reauthentication.
\end{itemize}

\chapter{Project Scope}
\section{Included Features}
The application includes the following user-facing areas:
\begin{itemize}[leftmargin=*]
    \item Authentication and profile screens.
    \item Classroom creation, joining, membership administration, and class content.
    \item Attendance sessions and attendance-sheet uploads.
    \item Notifications, class chat, and direct messaging.
    \item Media upload and file download support.
    \item An app-wide offline connectivity banner.
    \item Profile-based account deletion with typed confirmation and reauthentication.
\end{itemize}

\section{Out of Scope and Known Limitations}
The current application is an Android client backed by Firebase and Cloudinary; it
does not include a separate web client or a dedicated administrative portal. Account
deletion is coordinated by client-side Firestore operations rather than an atomic
server-side cascade. If cleanup fails partway through, some Firestore records may
already have been removed while the authentication account remains active. Files
uploaded to Cloudinary are not deleted by this flow. Owned classes require another
admin or member to receive ownership before deletion can proceed.

\chapter{Requirements}
\section{Functional Requirements}
\begin{enumerate}[leftmargin=*]
    \item Users can register and sign in with Firebase Authentication.
    \item Users can create a class or submit a join request using a class code.
    \item Class administrators can manage membership and classroom content.
    \item Users can access class announcements, materials, assignments, routines, and attendance.
    \item Users can communicate through class chat and direct conversations.
    \item The application displays an offline warning when Android reports no validated internet.
    \item Users can request account deletion after typing \texttt{DELETE} and providing their password.
\end{enumerate}

\section{Non-Functional Requirements}
\begin{itemize}[leftmargin=*]
    \item The UI should adapt to portrait and landscape device orientations.
    \item Authentication and Firestore operations should report failures to users.
    \item Firebase access must be governed by deployed Firestore Security Rules.
    \item Credentials and private API secrets must not be embedded in the Android app.
\end{itemize}

\chapter{System Design and Implementation}
\section{Technology Stack}
\begin{itemize}[leftmargin=*]
    \item \textbf{Language and UI:} Kotlin, Jetpack Compose, and Material 3.
    \item \textbf{Navigation:} Android Navigation Compose.
    \item \textbf{Authentication:} Firebase Authentication.
    \item \textbf{Application database:} Cloud Firestore.
    \item \textbf{Media:} Cloudinary Android SDK and Coil image loading.
    \item \textbf{Build:} Gradle with an Android application module.
\end{itemize}

\section{Architecture}
The application is a single-module Android project. The Compose screens manage UI
state and call Firebase services directly. \texttt{MainActivity} hosts navigation and
observes Android network connectivity to display the offline banner. Data classes
represent common Firestore records, while utility functions handle media transfers
and related helpers.

\section{Data Model}
The main Firestore collections include \texttt{users}, \texttt{classes},
\texttt{notifications}, and \texttt{directThreads}. Class-scoped subcollections
store posts, comments, materials, assignments, routines, attendance, chat messages,
join requests, and chat reports. Direct conversation messages are stored below each
thread document.

\section{Account Deletion Workflow}
The Profile screen asks the user to enter \texttt{DELETE} and re-enter the account
password. The client reauthenticates the email/password user, transfers owned
classrooms to an existing admin or member, removes supported linked Firestore data,
deletes the user's profile, and then deletes the Firebase Authentication account.
Deletion cannot proceed if a class has no eligible owner replacement. The matching
Firestore rules must be deployed for direct-thread participant removal.

\chapter{Classroom User Interface}
This page is reserved for the classroom screen from the running OMNIRA application.
Replace the placeholder below with a screenshot of the classroom interface. When
compiling from the repository root, a suggested image path is
\texttt{docs/images/classroom-ui.png}; create the directory and place the screenshot
there, or update the path below to match the image location.

\begin{figure}[H]
    \centering
    \IfFileExists{docs/images/classroom-ui.png}{
        \includegraphics[width=0.92\textwidth,height=0.72\textheight,keepaspectratio]
            {docs/images/classroom-ui.png}
    }{
        \fbox{
            \parbox[c][0.58\textheight][c]{0.88\textwidth}{
                \centering
                \Large Classroom UI Screenshot Slot\\[1em]
                \normalsize
                Add the classroom screenshot at
                \texttt{docs/images/classroom-ui.png}, then compile this report again.
            }
        }
    }
    \caption{OMNIRA classroom interface}
    \label{fig:classroom-ui}
\end{figure}

\chapter{Security and Production Considerations}
The application uses client-side Firebase SDK calls, so Firestore Security Rules are
essential: UI checks alone do not enforce access control. The repository rules preserve
broad authenticated access for some legacy collections and should be reviewed and
restricted before production use. The account-deletion workflow requires the updated
direct-message rules to be published. Client-side deletion is not atomic and should be
moved to a trusted backend if stronger consistency guarantees are required.

Cloudinary unsigned upload presets should be restricted. API secrets must remain out
of the app binary and source repository. Production deployments should also define
backup, retention, privacy, and incident-response procedures.

\chapter{Testing and Evaluation}
The Kotlin sources can be compiled with the repository's Gradle wrapper:

\begin{lstlisting}[language=bash]
./gradlew :app:compileDebugKotlin
\end{lstlisting}

Before release, test authentication, classroom membership, Firestore permissions,
offline and reconnect behavior, message operations, and account deletion using a
non-production Firebase project. Verify that class ownership transfers as intended
and that failures during cleanup are communicated clearly. Record test outcomes,
device models, Android versions, and any discovered defects here:

\vspace{1em}
\fbox{\parbox[c][4cm][c]{0.96\textwidth}{\centering
    \textit{Testing results and evaluation notes slot}
}}

\chapter{Conclusion and Future Work}
OMNIRA provides an integrated Android experience for classroom content, schedules,
attendance, notifications, and communication. The current implementation demonstrates
these workflows using Compose and Firebase services. Important future work includes
tightening Firestore authorization, adding automated tests for critical user journeys,
moving multi-step deletion to trusted server code, and supporting safe cleanup of
uploaded media.

\chapter*{References}
\addcontentsline{toc}{chapter}{References}
\begin{itemize}[leftmargin=*]
    \item Android Developers, \url{https://developer.android.com/}
    \item Firebase Documentation, \url{https://firebase.google.com/docs}
    \item Jetpack Compose Documentation, \url{https://developer.android.com/compose}
    \item Cloudinary Documentation, \url{https://cloudinary.com/documentation}
\end{itemize}

\end{document}
